\section*{Chapter \ref{ch:ml:build-an-order-book-model-end-to-end} --- Build: An Order-Book Model, End to End}

\begin{solution}{exo:ml:build-an-order-book-model-end-to-end:1}
Because production sees the market through the agent's own feed, with its latency, and computes features from what it has
received; research data recorded by the same agent with the same feature code are what production will see, and the offline and
online values can be tested equal. A full tape seen from nowhere would train the model on information, and timing, that the agent
never has.
\end{solution}

\begin{solution}{exo:ml:build-an-order-book-model-end-to-end:2}
Buying at the ask and selling at the bid a second later costs the spread, one tick, plus the taker fee twice, $2\times\$0.003=\$0.006$
a share, 0.6 ticks at one cent: 1.6 ticks a round trip. The mid moves by a tick or more within a second in about 5\% of cases, and
the forecasts are much smaller than a tick; no threshold leaves trades that cover the cost.
\end{solution}

\begin{solution}{exo:ml:build-an-order-book-model-end-to-end:3}
Without a model the quotes are filled 58.7 times a session with a one-second mark-out of 0.078 ticks: most fills come from orders
that move the price against the quote. The model withdraws the side about to be run over, so it is filled less (34.3 times) but
its fills are worth 0.405 ticks a second later. P\&L per session: \$21.73 against \$5.47.
\end{solution}

\begin{solution}{exo:ml:build-an-order-book-model-end-to-end:4}
Decisions are triggered at every book change, about four a second on average but many more in bursts. With $L=100$ milliseconds,
$\lambda L$ passes one in the bursts and decisions wait behind one another, so their median staleness (154 milliseconds) exceeds
the decision time. A production system conflates: when it finishes a decision it computes the next from the latest book, skipping
the updates it missed, so staleness is bounded by about twice the decision time; better still, it makes the decision fast enough
that $\lambda L$ stays well below one in the busiest bursts.
\end{solution}

\begin{solution}{exo:ml:build-an-order-book-model-end-to-end:5}
No: the model was profitable in all six sessions, and the monitor's large values come from the features that count events, which
move with the session's activity (2.8 to 6.4 book changes a second), while the price features stay between 0.02 and 0.13. The
monitor should use a reference spanning the range of activity (more training sessions, or history), normalise count features by
the session's activity, or watch counts with a separate activity monitor that pages the desk, not the \omterm{def:ml:the-machine-learning-team:roles}{model owner}.
\end{solution}

\begin{solution}{exo:ml:build-an-order-book-model-end-to-end:6}
Latency is worth money only where decisions compete with events or other traders. Here the P\&L is identical at 150 nanoseconds
and at 263 microseconds (\$21.73), and unchanged within its standard error up to 50 milliseconds: the speed-up buys nothing in
this market. It buys margin (a model that stays fast in bursts) and it buys everything in a market where competitors race in
microseconds.
\end{solution}

\begin{solution}{exo:ml:build-an-order-book-model-end-to-end:7}
The test sessions change their top of book about 4.3 times a second, and the P\&L crosses zero at about 106 milliseconds, some
0.45 book changes. At 2\,000 changes a second the same number of changes takes $106\times4.3/2000\approx0.23$ milliseconds, about
230 microseconds. The estimate leaves out competition (other fast traders take the stale quotes first, which moves the cliff to
the difference between their speed and the agent's), queue position (faster markets have longer queues), the different model a
busier market would need, and the burstiness that makes the decision queue grow.
\end{solution}

\begin{solution}{exo:ml:build-an-order-book-model-end-to-end:8}
Inputs: the four price features against a reference spanning many sessions, with a threshold calibrated to one false page a month
(chapter~27); the count features against the session's activity, owned by the desk. Outputs: the share of decisions with a side
withdrawn, the quoting uptime, the \omterm{def:ml:build-an-order-book-model-end-to-end:staleness}{decision staleness} distribution against a budget (median and 99th percentile, with an alarm when
$\lambda L$ in bursts approaches one). Outcomes: the one-second mark-out of fills and the session P\&L, with a CUSUM on the daily
mark-out, owned by the \omterm{def:ml:the-machine-learning-team:roles}{model owner}. Serving: parity of the C++ path with the reference on a daily sample of live inputs, owned by
the \omterm{def:ml:the-machine-learning-team:roles}{machine-learning engineer}.
\end{solution}

\begin{solution}{pb:ml:build-an-order-book-model-end-to-end:1}
\textbf{Part I.}
\begin{enumerate}
\item One venue of Book~10's simulator, one \$100 stock with a one-cent tick, rebate \$0.002 and taker fee \$0.003 a share,
\texttt{firm\_tape} background, 20-microsecond latencies; at every book change, ten features and a one-second forecast; quote one lot
at each allowed touch, withdrawing the side the forecast says will be run over ($\theta=0.1$ ticks), closing positions after 30
seconds; eight training, three validation and six test sessions of twenty minutes.
\item Touch imbalance, spread, order-flow imbalance over 1 and 5 seconds, signed volume over 1 and 5 seconds, trades over 5
seconds, mid change over 1 and 5 seconds, updates over 1 second; offline and online agree because the same agent and the same
\omterm{def:ml:data-and-feature-stores:store}{feature definitions} compute both.
\item The mid change in ticks over the next second, from the last event known at the decision; the mid usually does not move in a
second (89.4\% zeros).
\item Purged five-fold cross-validation with a one-second embargo; 7 leaves (0.558 against 0.555 and 0.548).
\end{enumerate}
\textbf{Part II.}
\begin{enumerate}[resume]
\item Validation information coefficient 0.522 for the forest, 0.452 for the TCN.
\item 153 nanoseconds at the median for the generated branches, against 263 microseconds for LightGBM's Python call and 166 for
the TCN.
\item The flat arrays, the C++ loop and branches and the Rust kernel reproduce LightGBM's predictions exactly on 200 shared vectors.
\item On the validation sessions at the compiled model's latency: \$53.9 per session at 0.1 ticks, the best of four.
\end{enumerate}
\textbf{Part III.}
\begin{enumerate}[resume]
\item \$21.73 per session (standard error \$3.90), positive in all six; quoting without a model earns \$5.47.
\item Model entries are worth 0.405 ticks one second later, against 0.078 without a model: the model declines the fills that
lose.
\item Largest indices from 0.14 to 1.60, all from count features that follow the session's activity; price features at most 0.13.
\item The time from the market event to the order's departure: the decision time, plus the wait behind earlier decisions when
they arrive faster than they can be made.
\end{enumerate}
\textbf{Part IV.}
\begin{enumerate}[resume]
\item \emph{Microseconds are money.} Per twenty-minute session, \$21.73 from zero to 263 microseconds, \$21.23 at 1 millisecond,
\$21.72 at 10, \$20.63 at 30, \$20.52 at 50, \$12.40 at 70, \$1.60 at 100 and $-\$12.27$ at 150; the one-second mark-out from
0.405 ticks to 0.148; the edge falls to zero at about 106 milliseconds, and below the no-model baseline at about 89.
\item Because the book changes about four times a second and nothing else races for the quotes: a withdrawal 50 milliseconds late
is rarely too late.
\item Staleness grows by a decision's own time and, in bursts, by the queue of decisions ($\lambda L>1$), exactly when the market
is busiest.
\item Far more book changes a second and fast competitors: the cliff would sit closer to zero by the ratio of event rates, and
closer still where competitors race.
\item Not for this market's sake; for robustness in bursts, and for the markets where it matters.
\item Competing fast agents, a reactive background (\texttt{firm.agentmkt}), realistic event rates, and conflation.
\item As in chapter~25 (the graph's stage keys as the code version, the forest as the artefact) and chapter~27 (the specification
of \cref{exo:ml:build-an-order-book-model-end-to-end:8}).
\item What the market's event rate and its competitors say it is: here nothing, elsewhere everything.
\end{enumerate}
\end{solution}

\begin{solution}{iq:ml:build-an-order-book-model-end-to-end:1}
Order-flow imbalance at the touch over short windows, the touch's size imbalance, the spread, signed traded volume, recent mid
changes and activity; built as stationary quantities over windows in event or clock time, as Cont, Kukanov and Stoikov and Kolm,
Turiel and Westray use.
\iqlookfor{order flow and imbalance, stationarity.}
\end{solution}

\begin{solution}{iq:ml:build-an-order-book-model-end-to-end:2}
Purge from the training folds every row whose label span overlaps the test fold, add an embargo after it, and keep sessions or days
of different roles on disjoint data; choose thresholds on data the model was not fitted on.
\iqlookfor{purging, embargo and disjoint test data.}
\end{solution}

\begin{solution}{iq:ml:build-an-order-book-model-end-to-end:3}
One \omterm{def:ml:data-and-feature-stores:store}{feature definition} executed by both paths (a \omterm{def:ml:data-and-feature-stores:store}{feature store}), research data recorded through the production code path, and a
test that offline values equal online ones at the same decision times, run in continuous integration and on samples in production.
\iqlookfor{shared code and a parity test.}
\end{solution}

\begin{solution}{iq:ml:build-an-order-book-model-end-to-end:4}
In the fills it declines: fewer fills with better mark-outs. Measure the mark-out curve of fills with and without the model and the
P\&L per session, on the same sessions.
\iqlookfor{adverse selection and mark-outs.}
\end{solution}

\begin{solution}{iq:ml:build-an-order-book-model-end-to-end:5}
Export the model to a flat form (tree arrays, quantised weights), generate or write a native kernel with no allocation on the hot
path, and check bit-for-bit or tolerance parity on shared test vectors; measure median and tail latency on the target machine.
\iqlookfor{compilation, parity vectors, measured tails.}
\end{solution}

\begin{solution}{iq:ml:build-an-order-book-model-end-to-end:6}
Trade the strategy at a grid of decision times in a simulator with realistic event rates and competitors, find where its P\&L
bends and where it crosses zero, and spend until the strategy's latency sits comfortably on the flat part, including in bursts.
\iqlookfor{a measured P\&L curve against latency.}
\end{solution}
